%1st paragraph
%
Demand for monitoring physical activity in office environments has increased.
\deleted{Typical monitoring targets include posture changes during seated deskwork, hand movements during PC operation, and movements to use office equipment or communal areas.}
\revised{Typical monitoring targets include office activities such as posture changes during seated work, hand movements during PC operation, and movements to use office equipment or communal areas.}
%
Physical activity monitoring using smart office technologies has become socially growing awareness, aiming to improve employees' working conditions and enhance overall organizational performance \cite{PAPAGIANNIDIS2020102027}.
%

%2nd paragraph
%
Wearable devices, smart furniture, and RGB cameras have been used to monitor physical movements, stress, and office occupancy \cite{6627764,mozos2017stress,ouwerkerk2013wireless,krejcar2019smart,9216892,vlaovic2022smart,ayers2001monitoring,antonino2019office}.
Light detection and ranging (LiDAR) acquires 3D data for recognizing indoor movements of people and objects \cite{9795869,oka2021spatial,10159661}.
\deleted{Multiple LiDAR sensors can cover indoor spaces without blind spots \cite{oka2021spatial,10159661}.}
\revised{Multiple LiDAR sensors provide the following benefits \cite{oka2021spatial,10159661}.
The sensors cover indoor spaces without blind spots.
Integration of multi-view point-cloud data increases point density.
Multi-view integration captures the full surface of a person.}

%3rd paragraph
Zenonos et al.\ presented a mood-prediction system using wearable physiological data \cite{7457166}.
Yang et al.\ presented a multi-camera indoor tracking system \cite{4284740}.
Behera et al.\ presented a multi-camera surveillance system for indoor environments \cite{6409058}.
Zenonos et al.\ relied on wearable physiological sensors \cite{7457166}.
Yang et al.\ and Behera et al.\ relied on RGB cameras for person detection and tracking \cite{4284740,6409058}.
These sensing modalities share the burden or privacy concerns described above.
The aforementioned studies have the following challenges:
studies using wearable devices have not considered the discomfort and burden associated with wearing such devices, and contain potential privacy risk.

%4th paragraph
% 提案システム
This paper proposes a system that estimates activity levels of office workers from spatial features extracted from 3D point-cloud data.
%
\deleted{Activity level is defined as the degree of concentration during PC-based deskwork.}
  %
  \deleted{The activity level indicates whether a person is focused on the task or has lost concentration.}
\deleted{Concentration is assumed to accompany active movements that appear as dynamic patterns in LiDAR data.}
\revised{Activity level is defined as a measure of learning efficiency during comprehension of meeting explanations and training content in offices.}
%
\revised{Numerical features represent 3D characteristics of point-cloud data, such as the spatial center of mass.}
%
Point-cloud data were collected during multiple short video viewing using multiple LiDAR sensors.
%
Activity indicators were obtained from objective and subjective evaluations of the video content.
%
Clustering was applied to the indicators to generate binary labels.
%
\deleted{Numerical features represent 3D characteristics of point-cloud data, such as the spatial center of mass.}
%
Binary labels were generated on the basis of the consistency of clustered activity indicators among dataset pairs.
%
Multiple models using different feature types were constructed, and their accuracy was evaluated using the binary labels.

%5th paragraph
\deleted{We previously presented an edge computing system that estimates activity levels of people working in office environments using a LiDAR sensor network \cite{kizawa2023estimation}.}
\revised{We previously presented an edge computing system that estimates physical activities of people in offices using a LiDAR sensor network \cite{kizawa2023estimation}.}
\deleted{Differently from the previous study, our study addresses estimation models using multiple features.
The previous study characterized the collected point-cloud data using a single method and a statistical model as the estimation model and using linear correlation as the evaluation metric.}
\revised{The differences between this study and the previous study are as follows.
The previous study estimated physical activities represented by a quantitative indicator from a single motion feature using a statistical model.
This study estimates activity levels represented by subjective and quantitative indicators from multiple features using machine-learning models.}
% A previous edge computing system used a single feature extraction method and a statistical model with linear correlation as the evaluation metric \cite{kizawa2023estimation}.
% The present study addresses estimation models using multiple features.
% The previous study extracted a single statistical feature from point-cloud frame similarity and evaluated estimation using linear correlation \cite{kizawa2023estimation}.

%6th paragraph
Section~\ref{sec:related-works} reviews related works.
Section~\ref{sec:proposed-system} describes the system model, feature extraction, and labeling.
Section~\ref{sec:evaluation} explains the experimental setup and machine-learning models.
Section~\ref{sec:evaluation-result} presents the evaluation results.
Section~\ref{sec:conclusion} concludes the paper.
